\sgaStableTarget{sga6.e13.sec.6}\subsection*{6. 附录: $P=\mathbf P^N_k$ 上 $(b)$-层的研究.}

\sgaStableTarget{sga6.e13.para.6-1}\textbf{6.1.} 设 $k$ 为代数闭域, $P=\mathbf P^N_k$. 设 $F$ 为 $P$ 上的凝聚层, 其 Hilbert 多项式为
\[
  \chi(F(n))=\sum_{i=0}^{r}a_i\binom{n+i}{i},
\]
又设
\[
  (c)=(c_0,\ldots,c_r)
\]
为满足 $(c)\geq(a)$ 的整数序列. 令
\[
  a_{m,i}=\sum_{j=0}^{r-i}a_{j+i}\binom{m-1+j}{j}
  \quad\text{及}\quad
  c_{m,i}=\sum_{j=0}^{r-i}c_{j+i}\binom{m-1+j}{j}.
\]
最后, 对每个整数 $q\geq1$, 用 $N_q$ 表示 $F$ 中支集维数至多为 $q-1$ 的最大凝聚子层, 并令 $F_q=F/N_q$.

\sgaStableTarget{sga6.e13.lem.6-2}\textbf{引理 6.2.} \emph{在 \hyperlink{sga6.e13.para.6-1}{6.1} 的条件下, 考虑正合列
\[
  0\longrightarrow I\longrightarrow\cO_P^{\oplus M}
  \longrightarrow F(m)\longrightarrow0
\]
其中 $m\geq0$. 则 $I$ 为 $p$-正则层, 这里
\(p=P_r(c_{m,0},\ldots,c_{m,r})\), $P_r$ 为 1.9 的第 $r$ 个
$(b)$-多项式.}

事实上, $\cO_P$ 显然为 $(0,\ldots,0,1)$-层 (共 $N+1$ 项), 所以由 1.6,
$I$ 是 $(0,\ldots,0,M)$-层. 由 \hyperlink{sga6.e13.rem.2-10}{2.10},
\(\chi(I(n))=M\binom{n+N}{N}
-\sum_{i=0}^{r}a_{m,i}\binom{n+i}{i}\). 断言由
\hyperlink{sga6.e13.thm.1-11}{1.11} 和 1.10 得出.

\sgaStableTarget{sga6.e13.prop.6-3}\textbf{命题 6.3.} \emph{在 \hyperlink{sga6.e13.para.6-1}{6.1} 的条件下, 假设 $F$ 为 $(b)$-层. 则可以利用关于 $a_i$、$b_i$ 和 $N$ 的泛多项式, 算出整数 $m,M,m_1,M_1$, 使存在正合列
\[
  \cO_P(-m_1)^{\oplus M_1}\longrightarrow
  \cO_P(-m)^{\oplus M}\longrightarrow F\longrightarrow0.
\]
此外, $N$ 只参与 $M_1$ 的计算.}

% [SGA6 current-rescribe idx667 | printed 654 | source-PDF 657 | high-res page 668]
事实上, 取 $m=P_{r-1}(c_0,\ldots,c_{r-1})$, 其中 $c_i=b_i-a_i$,
并取 $M=\sum_i a_i\binom{m+i}{i}$, 则由 \hyperlink{sga6.e13.thm.1-11}{1.11} 和 \hyperlink{sga6.e13.prop.1-3}{1.3}, 存在
正合列
\[
  0\longrightarrow I\longrightarrow\cO_P^{\oplus M}
  \longrightarrow F(m)\longrightarrow0.
\]
同样地, 考虑 \hyperlink{sga6.e13.lem.6-2}{6.2}, 取 $m_1=m+p$, 其中
$p=P_r(a_{m,0},\ldots,a_{m,r})$, 并取
$M_1=M\binom{p+N}{N}-\sum_i a_i\binom{m_1+i}{i}
\quad(=\chi(I(p)))$.

\sgaStableTarget{sga6.e13.thm.6-4}\textbf{定理 6.4.} \emph{在 \hyperlink{sga6.e13.para.6-1}{6.1} 的条件下, 假设
$F$ 是 $\cO_P(-m)^{\oplus M}$ 的商, 其中 $m\geq0$. 对 $i=0,\ldots,r$, 令
$b_i=P_{r-i}(c_{m,i},\ldots,c_{m,r})$, 这里 $P_s$
为 \hyperlink{sga6.e13.def.1-9}{1.9} 的第 $s$ 个 $(b)$-多项式; 令 $b=b_1+m-1$, 再令
$B=\sum_{i=1}^{r}a_i\binom{b+i}{i}$. 则:}
\begin{enumerate}[label=(\roman*)]
\item \emph{$F$ 为 $b$-正则层.}
\item \emph{$-B\leq a_0=h^0(F(b))-B$}
\item \emph{$F$ 为 $(b_0,\ldots,b_r)$-层.}
\end{enumerate}

事实上, 选取 $\cO_P(1)$ 的一个截面 $\sigma$, 可得如下行列均正合的交换图:
\[
\begin{tikzcd}[column sep=large,row sep=large]
&0 \arrow[d]&0 \arrow[d]&0 \arrow[d]&\\
0 \arrow[r]&I(-1) \arrow[r] \arrow[d]&I \arrow[r] \arrow[d]&J \arrow[r] \arrow[d]&0\\
0 \arrow[r]&\cO_P(-1)^{\oplus M} \arrow[r] \arrow[d]&\cO_P^{\oplus M} \arrow[r] \arrow[d]&\cO_Y^{\oplus M} \arrow[r] \arrow[d]&0\\
0 \arrow[r]&F(m-1) \arrow[r] \arrow[d]&F(m) \arrow[r] \arrow[d]&G(m) \arrow[r] \arrow[d]&0\\
&0&0&0&
\end{tikzcd}
\]

% [SGA6 current-rescribe idx668 | printed 655 | source-PDF 658 | high-res page 669]
由于根据 \hyperlink{sga6.e13.lem.1-7}{1.7} 有 $\chi(G(n))=\sum_{i=0}^{r-1}a_{i+1}\binom{n+i}{i}$,
由 \hyperlink{sga6.e13.lem.6-2}{6.2}, $J$ 为 $b_1$-正则层; 因而由 \hyperlink{sga6.e13.prop.1-4}{1.4(i)}, 当
$q\geq2$ 且 $p\geq b_1-q$ 时, 有 $H^q(I(p))=0$. 于是正合列
$H^q\bigl(\cO_P^{\oplus M}(p)\bigr)\longrightarrow
H^q(F(m+p))\longrightarrow H^{q+1}(I(p))$ 表明 $F(m)$ 为
$(b_1-1)$-正则层; 这就证明了 (i).

断言 (i) 蕴涵
$0\leq h^0(F(b))=\chi(F(b))
=\sum_{i=0}^{r}a_i\binom{b+i}{i}$, 这就证明了 (ii). 由 \hyperlink{sga6.e13.rem.2-10}{2.10} 和 \hyperlink{sga6.e13.def.1-9}{1.9},
\[
  \chi(F(b))
  =\sum_{i=0}^{r}a_{m,i}\binom{b_1-1+i}{i}
  \leq b_1+\sum_{i=0}^{r}c_{m,i}\binom{b_1-1+i}{i}
  =b_0.
\]
因此 $h^0(F(-1))\leq h^0(F(b))\leq b_0$. 对 $r$ 作归纳, 可设 $G$ 为 $(b_1,\ldots,b_r)$-层, 因而 (iii) 由 1.6.2 得出.

\sgaStableTarget{sga6.e13.lem.6-5}\textbf{引理 6.5.} \emph{在 \hyperlink{sga6.e13.para.6-1}{6.1} 的条件下, 假设 $F$
没有维数为 $0$ 的相伴素闭链.}
\begin{enumerate}[label=(\roman*)]
\item \emph{若 $h^0(F)\geq1$, 则 $h^0(F(-1))\leq h^0(F)-1$.}
\item \emph{当 $n\geq h^0(F)$ 时, 有 $H^0(F(-n))=0$.}
\item \emph{设 $0\to F(-1)\to F\to G\to0$ 正合, 并设
$n_0\geq0$ 满足 $H^0(G(-n_0))=0$. 则 $H^0(F(-n_0))=0$, 且
$h^0(F)\leq n_0\cdot h^0(G)$.}
\end{enumerate}

对于 (i), 设 $0\ne\sigma\in H^0(F)$. 截面 $\sigma$ 定义了 $F$ 的非零子层 $F'=\cO_X\cdot\sigma$, 所以 $s=\dim\Supp(F')\geq1$. 考虑如下行列均正合的交换图
\[
\begin{tikzcd}[column sep=large,row sep=large]
&0 \arrow[d]&0 \arrow[d]&0 \arrow[d]&\\
0 \arrow[r]&F'(-1) \arrow[r] \arrow[d]&F' \arrow[r] \arrow[d]&G' \arrow[r] \arrow[d]&0\\
0 \arrow[r]&F(-1) \arrow[r]&F \arrow[r]&G \arrow[r]&0.
\end{tikzcd}
\]

% [SGA6 current-rescribe idx669 | printed 656 | source-PDF 659 | high-res page 670]
由于 $s\geq1$, 有 $G'\ne0$; 因而 $F'$ 不包含于 $F(-1)$. 换言之, 截面 $\sigma$ 不属于单射 $H^0(F(-1))\to H^0(F)$ 的像, 这就证明了 (i).

断言 (ii) 由 (i) 立即得出, 只需注意到
$h^0(F)=0$ 显然蕴涵 $h^0(F(-1))=0$, 因为
$h^0(F(-1))\leq h^0(F)$. 在 (iii) 中得到
$0\leq h^0(F(n))-h^0(F(n-1))\leq h^0(G(n))$.
因此, 当 $n\geq n_0$ 时, 由于 $h^0(G(-n_0))=0$, 有 $h^0(F(-n_0))=h^0(F(-n))$; 所以由 (ii), $h^0(F(-n_0))=0$. 因而对 $p\geq1$,
$h^0(F(p-n_0))
\leq h^0(G(1-n_0))+\cdots+h^0(G(p-n_0))
\leq p\cdot h^0(G(p-n_0))$,
由此得到 (iii).

\sgaStableTarget{sga6.e13.prop.6-6}\textbf{命题 6.6.} \emph{在 \hyperlink{sga6.e13.para.6-1}{6.1} 的条件下, 假设
$F$ 为 $(b_0,\ldots,b_r)$-层. 则 $F_q$ 是一个
$\bigl((b_{q-1})^q,\ldots,(b_{q-1})^2,b_{q-1},b_q,\ldots,b_r\bigr)$-层.}

事实上, 设 $\sigma$ 为 $\cO_X(1)$ 的一个“一般”截面. 它定义如下行列均正合的交换图:
\[
\begin{tikzcd}[column sep=large,row sep=large]
&0 \arrow[d]&0 \arrow[d]&0 \arrow[d]&\\
0 \arrow[r]&N_q(-1) \arrow[r] \arrow[d]&N_q \arrow[r] \arrow[d]&G' \arrow[r] \arrow[d]&0\\
0 \arrow[r]&F(-1) \arrow[r] \arrow[d]&F \arrow[r] \arrow[d]&G \arrow[r] \arrow[d]&0\\
0 \arrow[r]&F_q(-1) \arrow[r] \arrow[d]&F_q \arrow[r] \arrow[d]&G'' \arrow[r] \arrow[d]&0\\
&0&0&0&
\end{tikzcd}
\]
其中由 \hyperlink{sga6.e13.prop.1-6}{1.6(i)}, $G$ 为 $(b_1,\ldots,b_r)$-层, 并且
$G''=G_{q-1}$ (这是深度的问题).

% [SGA6 current-rescribe idx670 | printed 657 | source-PDF 660 | high-res page 671]
先设 $q=1$. 则 $\dim\Supp(N_q)=0$, 所以
$H^1(N_q(-1))=0$, 且 $h^0(F_q(-1))\leq h^0(F(-1))\leq b_0$.
此外, $G'=0$, 从而 $G\simeq G''$. 因此 $F_1$ 确实为 $(b_0,\ldots,b_r)$-层. 又由\hyperlink{sga6.e13.lem.6-5}{ 6.5(ii)},
$H^0(F_1(-b_0))=0$. 若 $q\geq2$, 对 $q$ 作归纳. 可
设 $G_{q-1}$ 为
$\bigl((b_{q-1})^{q-1},\ldots,b_{q-1},\ldots,b_r\bigr)$-层, 且
$H^0(G_{q-1}(-b_{q-1}))=0$. 于是\hyperlink{sga6.e13.lem.6-5}{引理 6.5} 给出
$H^0(F_q(-b_{q-1}))=0$ 以及 $h^0(F_q(-1))\leq(b_{q-1})^q$,
断言得证.

\sgaStableTarget{sga6.e13.thm.6-7}\textbf{定理 6.7.} \emph{存在两个多项式序列
$\{A_i(X_0,\allowbreak\ldots,X_i;Y)\}$ 及
$\{A_i^{(q)}(X_0,\allowbreak\ldots,X_q;Y)\}$, 具有下述
性质. 假设满足 \hyperlink{sga6.e13.para.6-1}{6.1} 的条件, 且 $F$ 是
$\cO_P(-m)^{\oplus M}$ 的商, 其中 $m\geq0$.}
\begin{enumerate}[label=(\roman*)]
\item \emph{若 $F$ 为 $(b_0,\ldots,b_r)$-层, 则对
$i=0,\ldots,r$,}
\[
  |a_i|\leq A_{r-i}(b_i,\ldots,b_r;m).
\]
\item \emph{若
$\chi(F_q(n))=\sum_{i=0}^{r}a_i^{(q)}\binom{n+i}{i}$ 为 $F_q$ 的 Hilbert
多项式, 则对 $i=0,\ldots,q-1$,}
\[
  |a_i^{(q)}|\leq A_{r-i}^{(r-q)}(c_{q-1},\ldots,c_r;m).
\]
\end{enumerate}
\emph{(注意到 $a_{q-1}^{(q)}\leq a_{q-1}$, 且
$a_q^{(q)}=a_q,\ldots,a_r^{(q)}=a_r$, 因为
$\dim\Supp(N_q)\leq q-1$.)}

事实上, 取正合列
$0\longrightarrow F(-1)\longrightarrow F\longrightarrow G\longrightarrow0$, 使 $G$ 为 $(b_1,\ldots,b_r)$-层. 对 $r$ 作归纳, 并考虑 \hyperlink{sga6.e13.lem.1-7}{1.7}, 可
设 $A_0,\ldots,A_{r-1}$ 已经定义.
此时 \hyperlink{sga6.e13.thm.6-4}{6.4(ii)} 与 \hyperlink{sga6.e13.prop.1-8}{1.8} 一起给出
$A_r$ 的定义. 最后, 对 $F_q$ 代替 $F$ 应用 \hyperlink{sga6.e13.thm.6-4}{6.4(iii)}、\hyperlink{sga6.e13.prop.6-6}{6.6} 及断言
(i), 即在形式上得到 (ii).

% [SGA6 current-rescribe idx671 | printed 658 | source-PDF 661 | high-res page 672]
\sgaStableTarget{sga6.e13.cor.6-8}\textbf{推论 6.8. (Grothendieck).} \emph{设 $X$ 在
诺特概形 $S$ 上射影, $\cO_X(1)$ 对 $X/S$ 甚丰沛, 并设
$\mathcal J$ 为 $X/S$ 各纤维上凝聚层的同构类所成的一个族.
假设:}
\begin{enumerate}[label=(\alph*)]
\item \emph{$X$ 上存在凝聚层 $E$, 使 $\mathcal J$
包含于所有形如
$E_K$ 的层之商的同构类所成的族.}
\item[$(b_q)$] \emph{在 $F_K\in\mathcal J$ 的 Hilbert 多项式 $\chi(F_K(n))$ 中,
次数至少为 $q-1$ 的系数保持
有上界.}
\end{enumerate}
\emph{则 $(F_K)_q$ (当 $F_K$ 遍历 $\mathcal J$ 时) 构成
有界族; 此外, $\chi(F_K(n))$ 中次数
至少为 $q-2$ 的系数保持有下界.}

事实上, 显然可设 $X=\mathbf P^N_S$ 且
$E=\cO_X(-m)^{\oplus M}$. 第一个断言在形式上由\hyperlink{sga6.e13.thm.6-7}{
6.7(ii)} 及 1.13 得出. 第二个断言通过归纳, 利用正合列
$0\longrightarrow F(-1)\longrightarrow F\longrightarrow G\longrightarrow0$
及 \hyperlink{sga6.e13.thm.6-4}{6.4(ii)} 得出.

\sgaStableTarget{sga6.e13.def.6-9}\textbf{定义 6.9.} \emph{设 $k$ 为域. 所谓
\emph{$r$ 维特殊正 $k$-闭链}, 是指任意带有甚丰沛层 $\cO_X(1)$ 的射影
$k$-概形 $X$, 它是若干
$r$ 维闭子概形 $X_j$ 的并, 其中每个 $X_j$ 都是
一个整 $k_j$-概形 $X'_j$ (带有 $\cO_{X'_j}(1)$) 沿基变换 $\Spec(k)\to\Spec(k_j)$
的逆像. 所谓
$X$ 的\emph{次数}, 是指以下 Hilbert 多项式中的系数 $a_r$:
$\chi(\cO_X(n))=\sum_i a_i\binom{n+i}{i}$.}

\sgaStableTarget{sga6.e13.lem.6-10}\textbf{引理 6.10.} \emph{设 $k$ 为代数闭域, $X$ 为维数 $r$、次数 $d$ 的特殊正 $k$-闭链. 则 $\cO_X$ 为 $(0,\ldots,0,d)$-层.}

% [SGA6 current-rescribe idx672 | printed 659 | source-PDF 662 | high-res page 673]
事实上, 将 $X$ 视为 $\mathbf P^N_k$ 的子概形. 将 $k$ 换成某个纯超越扩张的代数闭包后 (由 \hyperlink{sga6.e13.prop.1-6}{1.6(i)} 允许如此), 可设 $k$ 含有在所有 $k_j$ 上代数无关的 $r(N+1)$ 个元素. 由这些元素定义的 $X$ 的连续 $r$ 个“一般”超平面截面, 便始终是次数 $d$ 的特殊正 $k$-闭链.

对 $r$ 作归纳, 只需注意, 当 $r=0$ 时,
有 $h^0(\cO_X(-1))=d$, 而当 $r\geq1$ 时, 有
$h^0(\cO_X(-1))=0$. 关于后一点, 单射
$\cO_X\longrightarrow\coprod_j\cO_{X_j}$ 使我们可设 $X$
本身来自一个整概形 $X'$. 然而此时
$h^0(\cO_X(-1))=h^0(\cO_{X'}(-1))=0$, 因为 $H^0(\cO_{X'})$ 为
域.

\sgaStableTarget{sga6.e13.cor.6-11}\textbf{推论 6.11.} (i) \emph{设 $r,d$ 为两个整数, 对 $i=0,\ldots,r$, 令
$c_i=A_{r-i}(0,\ldots,0,d;0)$ \hyperlink{sga6.e13.thm.6-7}{(6.7(i))}, 再令
$p=P_r(c_0,\ldots,c_r)$ \hyperlink{sga6.e13.def.1-9}{(1.9)}. 设 $k$ 为代数闭域,
$X$ 为维数至多 $r$、次数至多 $d$ 的特殊正 $k$-闭链,
其 Hilbert 多项式为
$\chi(\cO_X(n))=\sum_{i=0}^{r}a_i\binom{n+i}{i}$. 则对
$i=0,\ldots,r$,}
\[
  |a_i|\leq c_i.
\]
\emph{此外, $X$ 可嵌入为 $\mathbf P^N_k$ 的子概形, 其中
$N=d(r+1)-1$, 且由 (至多) $\binom{N+p}{N}$ 个次数为
$p$ 的方程定义.}

(ii) (Chow) \emph{设 $S$ 为 诺特概形. 当 $K$ 遍历
各 $s\in S$ 的剩余域 $k(s)$ 的代数闭扩张时,
维数和次数有界的特殊正 $K$-闭链构成一个
“有界族” (可将 $X$ 视为某个固定的
$\mathbf P^N_S$ 的子概形, 也可以在显然的抽象意义下理解).}

% [SGA6 current-rescribe idx673 | printed 660 | source-PDF 663 | high-res page 674]
事实上, 在 (i) 中, 第一个断言由 \hyperlink{sga6.e13.thm.6-7}{6.7(i)} 得出; 第二个断言
由 \hyperlink{sga6.e13.lem.6-10}{6.10}、\hyperlink{sga6.e13.prop.1-8}{1.8(i)}, 以及对
$0\longrightarrow I\longrightarrow\cO_P\longrightarrow\cO_X\longrightarrow0$,
$P=\mathbf P^N_k$ 应用 \hyperlink{sga6.e13.lem.6-2}{6.2}, 再结合 1.3 得出. 最后, (ii) 由 (i) 及
1.13.

\sgaStableTarget{sga6.e13.rem.6-12}\textbf{注 6.12.} 设 $k$ 为代数闭域,
$X$ 为纯 $r$ 维射影 $k$-概形, 没有嵌入的
相伴素闭链. 若允许 $\cO_X$ 中出现任意幂零元素,
则 Hilbert 多项式
$\chi(\cO_X(n))=\sum_{i=0}^{r}a_i\binom{n+i}{i}$ 中的系数, 就不能再仅用
$a_r$ 来控制其界.

例如, 设 $Z$ 为 $\mathbf P^3_k$ 中次数 $d$ 的光滑曲线.
对每个 $n\gg0$, 存在一个包含 $Z$ 的次数 $n$ 的光滑曲面 $Y$.
于是 $Z$ 是 $Y$ 上的 Cartier 除子; 令
$\cO_{X_n}=\cO_Y/\cO_Y(-2Z)$. 有正合列
$0\longrightarrow\cO_X(-Z)\longrightarrow
\cO_X\longrightarrow\cO_Z\longrightarrow0$\footnote{译校注: 原书第660页此列两处下标均印作 $X$, 本段只定义了 $X_n$. 对二重除子 $X_n=2Z$, 通常所用的正合列为 $0\to\cO_Z(-Z)\to\cO_{X_n}\to\cO_Z\to0$. 此处保留原印, 不以这一解释替换原文.}, 以及公式
$\langle Z^2\rangle_Y=2p_a(Z)-2-(n-4)d$. 由此得出, 当 $n\longrightarrow\infty$ 时,
$a_0=\chi(\cO_{X_n}(-1))\longrightarrow\infty$,
但 $a_1=2d$.

\sgaStableTarget{sga6.e13.cor.6-13}\textbf{推论 6.13.} \emph{在 \hyperlink{sga6.e13.para.6-1}{6.1} 的条件下, 考虑正合列
\[
  \cO_P(-m_1)^{\oplus M_1}\longrightarrow
  \cO_P(-m)^{\oplus M}\xrightarrow{\alpha}F\longrightarrow0
\]
其中 $m_1\geq m\geq0$. 则对 $i=0,\ldots,N$,
\[
% SGA-SRC-001438
  |a_i|\leq A_{N-i}(0,\ldots,0,M;m_1)
  +M\binom{m}{N-i}.
\]}

事实上, 设 $I=\Ker(\alpha)$. 显然, $I$ 是
$(0,\ldots,0,M)$-层, 且
$\chi(I(n))=\allowbreak\sum_{i=0}^{N}
\bigl[M\tbinom{-m-1+N-i}{N-i}-a_i\bigr]\allowbreak\tbinom{n+i}{i}$.
断言由 \hyperlink{sga6.e13.thm.6-7}{6.7(i)} 及公式
$\tbinom{-m-1+Q}{Q}=(-1)^Q\tbinom{m}{Q}$.

% [SGA6 current-rescribe idx674 | printed 661 | source-PDF 664 | high-res page 675]
\sgaStableTarget{sga6.e13.cor.6-14}\textbf{推论 6.14. (I. Hermann).} \emph{对每个 $N\geq0$, 存在
一个多项式 $R_N(m)$, 使对每个域 $k_0$ 及多项式环 $k_0[T_1,\ldots,T_N]$
中每个由次数至多为 $m$ 的元素生成的理想
$I$, 其根理想 $\sqrt I$ 都可由次数至多为
$R_N(m)$ 的元素生成.}

事实上, 设 $k$ 为 $k_0$ 的代数闭包. 引入一个
辅助不定元后, 用 $Y$ (相应地, $X$) 表示由 $I$
(相应地, $\sqrt I$) 的齐次化理想定义的 $P=\mathbf P^N_k$ 的闭
子概形. 在 Hilbert 多项式
$\chi(\cO_Y(n))=\sum_{i=0}^{N}b_i\binom{n+i}{i}$ 中, 由 6.13, 有
$|b_i|\leq c_i$, 其中 $c_i=A_{N-i}(0,\ldots,0,1;m)$.

设 $X=\bigcup X^q$ 为按维数作的分解. 则有
$\deg(X^q)\leq e_q$, 其中 $e_q=P_{N-q}(c_q,\ldots,c_N)$.
事实上, 用一个余维 $q$ 的“一般”线性子空间截 $X$ (相应地, $Y$) 后, 可设 $q=0$. 此时
$\deg(X^q)=h^0(\cO_{X^q})$ 且
$h^0(\cO_{X^q})\leq h^0(\cO_Y)$, 而 \hyperlink{sga6.e13.thm.6-4}{6.4(iii)} 给出
$h^0(\cO_Y)\leq e_q$.

因此, 由 6.10, $\cO_{X^q}$ 为 $(0,\ldots,0,e_q)$-层. 由
单射 $\cO_X\longrightarrow\prod_q\cO_{X^q}$ 可知,
$\cO_X$ 为 $(e_0,\ldots,e_N)$-层. 因而在 Hilbert 多项式
$\chi(\cO_X(n))=\sum_{i=0}^{n}a_i\binom{n+i}{i}$\footnote{译校注: 原书第661页求和上限印作 $n$; 按周围的 $N$ 维射影空间及系数序列, 自然应为 $N$. 本文保留印本.} 中, 有
$|a_i|\leq f_i$, 其中由 6.7, $f_i=A_{N-i}(e_i,\ldots,e_N;0)$.
最后, $X$ 的理想 $\sqrt I$ 为 $R_N(m)$-正则的, 其中
$R_N(m)=P_N(f_0,\ldots,f_N)$, 断言由 1.3 得出.

% [SGA6 current-rescribe idx675 | printed 662 | source-PDF 665 | high-res page 676]
\sgaStableTarget{sga6.e13.sec.7}\subsection*{7. 附录: Hodge 指标定理}

这里通过直接论证, 重新得到并补充前面的一些结果.

\sgaStableTarget{sga6.e13.thm.7-1}\textbf{定理 7.1.} \emph{设 $X$ 为代数闭域上的固有不可约
曲面. 假设 $X$ 上存在可逆
层 $H$, 使 $\langle c_1(H)^2\rangle>0$ (例如,
$H$ 丰沛). 则 $X$ 上满足
\[
  \langle c_1(L)\cdot c_1(H)\rangle=0,
  \qquad
  \langle c_1(L)^2\rangle\geq0,
\]
的可逆层 $L$, 都与 $0$ 数值等价. 因此, 由 \hyperlink{sga6.e13.thm.4-6}{4.6}, 相交型在 $A^1_{\mathrm{num}}(X)_{\Q}
=\bigl[\Pic(X)/\Pic^\tau(X)\bigr]\otimes\Q$ 上的二次型非退化, 其型为
$(1,\rho-1)$.}

\sgaStableTarget{sga6.e13.para.7-1-1}\textbf{7.1.1.} 先假设 $X$ 既约 (因而整),
$H$ 丰沛, $\langle c_1(L)\cdot c_1(H)\rangle=0$, 且
$\langle c_1(L)^2\rangle>0$. 用反证法, 选取 $m>0$,
使 $H_1=L\otimes H^{\otimes m}$ 丰沛, 并选取 $n$,
使令 $N=L^{\otimes n}\otimes H^{-1}$ 后, 有
$\langle c_1(N)\cdot c_1(H_1)\rangle
=n\langle c_1(L)^2\rangle-m\langle c_1(H)^2\rangle>0$. 还有
$\langle c_1(N)\cdot c_1(H)\rangle=-\langle c_1(H)^2\rangle<0$ 及
$\langle c_1(N)^2\rangle
=n^2\langle c_1(L)^2\rangle+\langle c_1(H)^2\rangle>0$, 这与
下述引理矛盾.

\sgaStableTarget{sga6.e13.lem.7-1-2}\textbf{引理 7.1.2.} \emph{设 $X$ 为代数闭域上的射影整曲面. 对 $X$ 上满足 $\langle c_1(N)^2\rangle>0$ 的可逆层 $N$, 下列条件等价:
\begin{enumerate}[label=(\roman*)]
\item 对每个丰沛层 $H$, 有 $\langle c_1(N)\cdot c_1(H)\rangle>0$;
\item[(i$'$)] 存在丰沛层 $H$, 使 $\langle c_1(N)\cdot c_1(H)\rangle>0$;
\item 存在整数 $n>0$, 使 $H^0(X,N^{\otimes n})\ne0$.
\end{enumerate}}

% [SGA6 current-rescribe idx676 | printed 663 | source-PDF 666 | high-res page 677]
事实上, (ii) 蕴涵存在有效除子 $D$, 使
$N^{\otimes n}\simeq\cO_X(D)$, 而 $\langle c_1(N)^2\rangle>0$ 蕴涵
$D>0$; 因而得到 (i).

为证明 (i$'$) 蕴涵 (ii), 显然可设
$H=\cO_X(1)$ 甚丰沛. 于是存在闭的整局部
主曲线 $Y$, 使 $H=\cO_X(Y)$. 由 Riemann 定理,
(i$'$) 蕴涵当 $p\geq0$ 及
$n\geq n_0$ 时, 有 $H^1\bigl(N_Y^{\otimes n}(p)\bigr)=0$, 其中 $n_0=2h^1(\cO_Y)-1$.
由于当 $p\gg0$ 时 $H^2(N^{\otimes n}(p))=0$, 正合列
\[
  H^1\bigl(N_Y^{\otimes n}(p)\bigr)\longrightarrow
  H^2\bigl(N^{\otimes n}(p-1)\bigr)\longrightarrow
  H^2\bigl(N^{\otimes n}(p)\bigr)\longrightarrow0
\]
给出当 $n\geq n_0$ 时 $H^2(N^{\otimes n})=0$. 因此, 当 $n\gg0$ 时,
\[
  h^0(N^{\otimes n})
  =\chi(N^{\otimes n})+h^1(N^{\otimes n})
  \geq\chi(N^{\otimes n})
  =\frac12\langle c_1(N)^2\rangle n^2+\cdots>0,
\]
引理证明完毕.

\sgaStableTarget{sga6.e13.para.7-1-3}\textbf{7.1.3.} 现设在 \hyperlink{sga6.e13.thm.7-1}{7.1} 中, $X$ 上存在可逆层 $M$,
使 $\langle c_1(L)\cdot c_1(M)\rangle\ne0$, 并用
反证法. 把 $X$ 换成 $X$ 的一个 Chow
表示的底层既约概形; 由投影公式 \hyperlink{sga6.e10.eq.4-4-4}{X
4.4.4} 允许这样做. 设 $G$ 在 $X$ 上丰沛. 选取 $m>0$, 使
$H_1=G^{\otimes m}\otimes M$ 丰沛, 且
$\langle c_1(L)\cdot c_1(H_1)\rangle\ne0$. 选取 $p,q$, 其中 $q\ne0$,
使令 $L_1=L^{\otimes p}\otimes H^{\otimes q}$ 后, 有
$\langle c_1(L_1)\cdot c_1(H_1)\rangle
=p\langle c_1(L)\cdot c_1(H_1)\rangle
+q\langle c_1(H)\cdot c_1(H_1)\rangle=0$. 此时还有
$\langle c_1(L_1)^2\rangle=p^2\langle c_1(L)^2\rangle
+q^2\langle c_1(H)^2\rangle>0$. 这就是 \hyperlink{sga6.e13.para.7-1-1}{7.1.1} 的情形, 其中 $L=L_1$
及 $H=H_1$.

\sgaStableTarget{sga6.e13.para.7-1-4}\textbf{7.1.4.} 最后考虑 7.1 的一般情形. 仍用反证法, 假设 $L$ 不与 $0$ 数值等价. 把 $X$ 换成 $X_{\red}$ 的正规化, 于是 $X$

% [SGA6 current-rescribe idx677 | printed 664 | source-PDF 667 | high-res page 678]
只有有限多个奇点; 这一归结由投影公式 \hyperlink{sga6.e10.eq.4-4-4}{X 4.4.4} 及 \hyperlink{sga6.e13.prop.4-5}{4.5(ii)} 保证. 因而 $X$ 上存在闭整曲线 $Y$, 使
$\langle c_1(L)\cdot Y\rangle\ne0$, 且 $Y$ 在一个有限集合 $Z$
之外局部为主. 设 $f:X'\to X$ 为沿定义 $Y$ 的理想
$I\subset\cO_Y$\footnote{译校注: 原书第664页将 $I$ 印为 $\cO_Y$ 的理想; 要在 $X$ 中定义 $Y$ 并沿它爆破, 自然应为 $\cO_X$ 的理想. 本文保留印本.} 所作的爆破. 理想 $\mathcal M=I\cO_{X'}$
可逆, 并定义 $X'$ 的一个子概形, 它是两个
子概形 $Y'$ 与 $Y''$ 的并, 满足 $f(Y')=Y$ 及 $f(Y'')\subset Z$.
由投影公式 \hyperlink{sga6.e10.eq.4-4-4}{X 4.4.4},
$\langle c_1(f^*L)\cdot c_1(\mathcal M)\rangle
=-\langle c_1(f^*L)\cdot Y'\rangle
-\langle c_1(f^*L)\cdot Y''\rangle
=-\langle c_1(L)\cdot Y\rangle-0\ne0$. 因而可将 $X,L,H$
换成 $X',f^*L,f^*H$. 这就归结到前一种情形,
并完成 7.1 的证明.

\sgaStableTarget{sga6.e13.rem.7-2}\textbf{注 7.2.} 尚不知道, 若代数闭域上的不可约固有曲面具有可逆层 $H$, 满足 $\langle c_1(H)^2\rangle>0$, 它是否必为射影曲面. 也不知道对每个不可约固有曲面, $A^1_{\mathrm{num}}(X)_{\Q}$ 上的相交型是否非退化. 由 \hyperlink{sga6.e13.thm.7-1}{7.1} 可知, 若它退化, 则必为负的 (非定的).

% SGA-SRC-001442 — reconstructed by mathematical inference; scan inconclusive
\sgaStableTarget{sga6.e13.cor.7-3}\textbf{推论 7.3.} \emph{在 \hyperlink{sga6.e13.thm.7-1}{7.1} 的条件下, 设 $M$ 为
$X$ 上的可逆层, 它与 $H$ 不数值等价\footnote{译校注: 原书第664页的假设为“不数值等价”, 但严格不等式需要 $M$ 与 $H$ 在数值上不成比例; 例如 $M=H^{\otimes2}$ 满足原假设却使所列两端相等. 本文保留原陈述.},
并令 $d=\langle c_1(M)\cdot c_1(H)\rangle$ 及
$h=\langle c_1(H)^2\rangle$. 则
$\langle c_1(M)^2\rangle<d^2/h$.}

% SGA-SRC-001441
事实上, 令 $L=M^{\otimes h}\otimes H^{\otimes(-d)}$. 由于
$\langle c_1(L)\cdot c_1(H)\rangle=0$, \hyperlink{sga6.e13.thm.7-1}{定理 7.1} 给出
$\langle c_1(L)^2\rangle<0$, 断言得证.

% [SGA6 current-rescribe idx678 | printed 665 | source-PDF 668 | high-res page 679]
\sgaStableTarget{sga6.e13.cor.7-4}\textbf{推论 7.4.} \emph{设 $X$ 为代数闭域上的射影不可约概形,
维数 $r\geq2$, 并
设 $H$ 为 $X$ 上的丰沛层. 则:}
\begin{enumerate}[label=(\roman*)]
\item \emph{$X$ 上的可逆层 $L$ 与
$0$ 数值等价, 当且仅当}
\[
  \langle c_1(L)\cdot c_1(H)^{r-1}\rangle=0,
  \qquad
  \langle c_1(L)^2\cdot c_1(H)^{r-2}\rangle\geq0;
\]
\item \emph{对 $X$ 上每个不与 $H$
数值等价的可逆层 $M$, 有\footnote{译校注: 原书第665页沿用“不数值等价”的假设. 与前页推论7.3相同, 严格不等式自然需要数值上不成比例; 取 $M=H^{\otimes2}$ 时两端相等. 此处保留原陈述.}}
\[
  \langle c_1(M)^2\cdot c_1(H)^{r-2}\rangle
  <\frac{\langle c_1(M)\cdot c_1(H)^{r-1}\rangle^2}
  {\langle c_1(H)^r\rangle}.
\]
\end{enumerate}

断言 (ii) 由 (i) 用 7.3 的论证立即得出. 为
证明 (i), 设 $Y$ 为 $X$ 中的闭整曲线; 须证明
$\langle c_1(L)\cdot Y\rangle=0$. 用 $H^{\otimes n}$ 将 $X$ 嵌入 $\mathbf P^N_k$,
其中 $n>0$. 设 $I=\bigoplus I_\nu$ 为 $Y$ 的
齐次理想, 取整数 $\nu_0$, 使
当 $\nu\geq\nu_0$ 时 $I_{\nu_0}$ 生成 $I_\nu$. 于是,
包含 $Y$ 的次数 $\nu_0$ 的超曲面 $Z$,
在 $Y$ 外没有公共点. 对 $r=\dim X$ 作归纳, 利用 Bertini 定理, 可找到
包含 $Y$ 的 $r-2$ 个次数 $\nu_0$ 的超曲面 $Z_1,\ldots,Z_{r-2}$,
使 $X'=Z_1\cap\cdots\cap Z_{r-2}\cap X$ 不可约.
由于 $\langle c_1(L)\cdot Y\rangle
=\langle c_1(L|_{X'})\cdot Y\rangle$, 只需证明
$L|_{X'}$ 在 $X'$ 上与 $0$ 数值等价. 由投影
公式, $\langle c_1(L|_{X'})\cdot c_1(H|_{X'})\rangle=0$, 且
$\langle c_1(L|_{X'})^2\rangle\geq0$, 因而断言由 7.1 得出.

\sgaStableTarget{sga6.e13.cor.7-5}\textbf{推论 7.5.} \emph{设 $k$ 为代数闭域,
$Y$ 为射影 $k$-概形, 带有丰沛层 $H$, 并
设 $X$ 为 $H$ 的截面 $\varphi$ 的零点概形. 假设
对 $Y$ 的每个不可约分支 $Y^i$,
$X\cap Y^i=V(\varphi|_{Y^i})$ 的不可约分支
% [SGA6 current-rescribe idx679 | printed 666 | source-PDF 669 | high-res page 680]
的维数至少为 $2$. 则态射
$A^1_{\mathrm{num}}(Y)\longrightarrow A^1_{\mathrm{num}}(X)$ 为
单射.}

事实上, 显然可设 $Y$ 为整概形, 断言由 \hyperlink{sga6.e13.cor.7-4}{7.4(i)} 及投影公式得出.

\sgaStableTarget{sga6.e13.bibliography}\subsection*{参考文献}
\begin{enumerate}[label={[\arabic*]}]
\item S. Abhyankar, ``Resolution of singularities of arithmetical surfaces,'' 载于 \emph{Arithmetical Algebraic Geometry}, Purdue University, Harper \& Row, New York, 1965.
\item D. Mumford, ``Pathologies III,'' \emph{Amer. J. Math.}, 1967, pp. 94--103. \emph{Amer. J. Math.}, 1967, pp. 94--103.
\item A. Grothendieck 与 J. Dieudonn\'e, \emph{EGA, \'El\'ements de G\'eom\'etrie Alg\'ebrique}, Publ. Math. Inst. des Hautes \'Etudes Sci., 1960.
\item A. Grothendieck, \emph{FGA, Fondements de la g\'eom\'etrie alg\'ebrique}, S\'eminaire Bourbaki 1957--1962 的摘录.
% SGA-SRC-001447
\item A. Grothendieck et al., \emph{SGA, S\'eminaire de g\'eom\'etrie alg\'ebrique}, 1960--66, I.H.E.S.
\item O. Jussila, “Formalisme des intersections sur les sch\'emas alg\'ebriques propres”, 本讨论班第 X 讲.
\item M. Raynaud, “Un th\'eor\`eme de repr\'esentabilit\'e relative sur le foncteur de Picard”, 本讨论班第 XII 讲.
\end{enumerate}

% [SGA6 current-rescribe idx680 | printed 667 | source-PDF 670 | high-res page 681]
